\documentclass[11pt,a4paper]{article}
\usepackage[T1]{fontenc}
\usepackage{mathpazo}
\usepackage[a4paper,margin=2.5cm]{geometry}
\pagestyle{empty}
\setlength{\parskip}{8pt}
\setlength{\parindent}{0pt}
\linespread{1.1}

\begin{document}

{\large\bfseries Measuring the External Field: A Cortical-Proxy Content Observable and the
Mean-Field Bound on Media-Driven Opinion Change}

\vspace{4pt}
Brian Mwai, 1050555, B.Sc.\ Physics, Department of Natural Sciences, The Catholic University
of Eastern Africa. Supervisor: Dr.\ Songa Mutambi.

\vspace{2pt}\hrule\vspace{8pt}

\textbf{Abstract}

Physicists often model how opinions spread through a society with the Ising model, the same
mathematics that describes how tiny magnets in a metal line up with each other. In these
models, news and other media act as an outside push on the population. The strength of that
push is always assumed and never measured. This project asks whether it can be measured from
the content itself.

Each news article was converted to speech and passed through TRIBE v2, a published AI model
from Meta that predicts how the human brain would respond to it. From that prediction, the
project compares activity in brain regions linked to emotion with activity in regions linked
to careful reasoning, which gives one number per article. No person was scanned: every value
is a prediction made by the model.

Four hundred articles in four groups (fear-driven, outrage, clickbait and neutral) were
measured twice. Fear-driven articles scored clearly higher than neutral ones, and the result
repeated on the second run. However, each group came from a different source, so the
difference cannot be credited to the writing style alone. A simple word-counting method also
sorted manipulative from neutral articles better than this measure, so the measure is not a
detector of manipulation.

The main physics result is a formula for the smallest push media must have before it can
flip a population's majority opinion in this type of model. Using the spread of scores
measured here, the formula sets a minimum that any future model of media influence must meet.

\textbf{Keywords:} opinion dynamics, Ising model, media influence, brain-response prediction.

\end{document}
